% ===== s04-bung.tex —— 原文 §2（p.16–p.23）=============================
% 底本：Laurent Fargues, Geometrization of the local Langlands correspondence:
%       an overview, 2016-02-02，§2. The stack Bun_G（p.16 起，p.23 §3 前止）。
% 说明：原文 "Hope 2.3" 用 hope 环境（已由 Lead 加入 preamble.tex:78，
%       与 theorem 共用计数器，输出「期望」）。
\section{叠 \texorpdfstring{$\Bun_G$}{Bun\_G}}% 原文 2. The stack Bun_G
\srcpage{16}

\subsection{定义与初步性质}% 原文 2.1. Definition and first properties

设 $E$ 或者是 $\FF_q((\pi))$，或者是 $\QQ_p$ 的有限次扩域、其剩余域为 $\FF_q$。设 $G$ 是 $E$ 上的约化群。考虑 $\mathrm{Perf}_{\FF_q}$，\emph{配以 pro-étale 拓扑}。

\begin{definition}\label{def:2.1}
我们记 $\Bun_G$ 为 $\mathrm{Perf}_{\FF_q}$ 上的纤维范畴
\[
S\longmapsto \text{群胚}\bigl(\text{从 } \Rep(G) \text{ 到 } \Bun_{X_S} \text{ 的 } \otimes\text{-正合函子}\bigr).
\]
\end{definition}

\yizhu{$\otimes$-正合函子指与张量积相容的正合函子；按 Tannaka 观点，$G$-丛由 $\Rep(G)$ 上的纤维函子刻画。}{A：原文 §2.1 p.16}

下面的技术性结果，读者可以直接承认。

\begin{proposition}\label{prop:2.2}
$\Bun_G$ 是 $\mathrm{Perf}_{\FF_q}$ 上的叠。
\end{proposition}

\begin{hope}\label{hope:2.3}
$\Bun_G$ 是一个“光滑菱形叠”。
\end{hope}

\yizhu{原文用 Hope（期望），而非后面的 Conjecture（猜想）；引号内的“光滑”当时尚无定义，作者在下段明确承认这一点。}{A：原文 §2.1 p.16}

\emph{迄今为止，这个概念还没有精确的定义。尽管如此，下面给出几个事实，试图说服读者。我们处理 $\mathrm{GL}_n$ 的情形。用标准技巧可以化归到这一情形。事实上，若固定一个嵌入 $G\subset \mathrm{GL}_n$，则存在线性表示 $\rho:\mathrm{GL}_n\to \mathrm{GL}(W)$ 以及 $W$ 中的一条直线 $D$，使得 $G$ 是 $D$ 在 $\mathrm{GL}_n$ 中的稳定化子。此时，$G$-丛等同于一个秩 $n$ 的向量丛 $\mathcal{E}$，连同一个 $\rho_*\mathcal{E}$ 的线子丛，该子丛在局部是直和项。}

\begin{proposition}\label{prop:2.4}
设 $S\in \mathrm{Perf}_{\FF_q}$，$\mathcal{E}_1,\mathcal{E}_2$ 是 $X_S$ 上两个常数秩相同的向量丛。则层
\[
T/S\longmapsto \Isom\bigl(\mathcal{E}_1|_{X_T},\mathcal{E}_2|_{X_T}\bigr)
\]
是 $S$ 上的菱形。当 $E=\FF_q((\pi))$ 时，它可由 $S$ 上的完满空间表示。
\end{proposition}

事实上，$\Isom\subset \mathrm{Hom}$ 相对可表为开子集。为此，设 $u:\mathcal{E}_1|_{X_T}\to \mathcal{E}_2|_{X_T}$ 是态射，并考虑使 $u$ 成为同构的最大开子集 $U\subset X_T$。注意 $f:|X_T|\to |T|$ 是连续闭映射。则 $T\setminus f(|X_T|\setminus U)\subset T$ 是开子集，使得 $T'\to T$ 经由 $U$ 分解，当且仅当 $u|_{X_{T'}}$ 是同构。

因此只需证明 $T/S\mapsto \Gamma\bigl(X_T,(\mathcal{E}_1^\vee\otimes \mathcal{E}_2)|_{X_T}\bigr)$ 可表。这是下面结果的一个推论。

\begin{proposition}\label{prop:2.5}
设 $\mathcal{E}$ 是 $X_S$ 上的向量丛，$f:(X_S)_{\text{pro-ét}}\to S_{\text{pro-ét}}$。则 $f_*\mathcal{E}$ 可由菱形表示；当 $E=\FF_q((\pi))$ 时，可由完满空间表示。
\end{proposition}

可以假定 $S$ 是 affinoid 完满的。此时 Kedlaya 与 Liu 的一个定理（[16]）说：当 $d\gg0$ 时，$\mathcal{E}(d)$ 由其整体截面生成（也就是说 $\mathcal{O}(1)$ 是丰沛的）。把它用于 $\mathcal{E}^\vee$，可知存在 $d_1,d_2\in\ZZ$、$n_1,n_2\in\mathbb{N}$ 以及一个部分消解
\[
0\longrightarrow \mathcal{E}\longrightarrow \mathcal{O}(d_1)^{n_1}\longrightarrow \mathcal{O}(d_2)^{n_2}.
\]
这给出部分消解
\[
0\longrightarrow f_*\mathcal{E}\longrightarrow f_*\mathcal{O}(d_1)^{n_1}\longrightarrow f_*\mathcal{O}(d_2)^{n_2}.
\]
但现在，对一切 $d\in\ZZ$，$f_*\mathcal{O}(d)$ 可由菱形表示；当 $E=\FF_q((\pi))$ 时，可由完满空间表示。对 $E|\QQ_p$ 的情形，这是 Banach–Colmez 空间理论的结果（[4]）。当 $E=\FF_q((\pi))$ 时，若 $d<0$ 它是 $0$；若 $d=0$ 它是 $E$；若 $d>0$ 它由 $S$ 上 $d$ 维的完满开球表示。更确切地说，对 $E=\FF_q((\pi))$ 的情形：\srcpage{17}

\begin{itemize}
\item 若 $S=\Spa(R,R^+)$，则 $\mathcal{E}|_S=\bigcup_{n\ge0}\Spa\bigl(\mathcal{C}(\pi^{-n}\mathcal{O}_E,R),\mathcal{C}(\pi^{-n}\mathcal{O}_E,R^+)\bigr)$。
\item 若 $d>0$，用 1.5.2 节的记号，$f_*\mathcal{O}(d)=\bigl(\varprojlim_{\mathrm{Frob}}\mathcal{G}_d\bigr)\times S$，其中 $\varprojlim_{\mathrm{Frob}}\mathcal{G}_d=\varprojlim_{\times\pi}\mathcal{G}_d$ 是形式群 $\mathcal{G}_d$（一个形式 $E$-向量空间）的万有覆盖。
\end{itemize}

\yizhu{$d$ 维完满开球指同构于 $\Spa\bigl(R\langle T_1^{1/p^\infty},\dots,T_d^{1/p^\infty}\rangle\bigr)$ 的完满空间；例 2.7 中 $d=1$ 的情形即 $\Spf\bigl(\FF_q[\![x^{1/p^\infty}]\!]\bigr)$。}{A：原文 §2.1 p.16–18 及例 2.7}

于是 $f_*\mathcal{E}$ 是群菱形（相应地，完满空间）之间态射的核，因而可表。

\begin{corollary}\label{cor:2.6}
$\Bun_G$ 的对角线可由菱形表示；当 $E=\FF_q((\pi))$ 时，可由完满空间表示。
\end{corollary}

现在的问题是：$\Bun_G$ 是否存在“光滑”表示。首先注意，$\Spa(\FF_q((T^{1/p^\infty})))\to \Spa(\FF_q)$ 是解析拓扑下的满态射。事实上，若 $(R,R^+)$ 是 affinoid 完满的，则每选一个伪一致化元 $\varpi_R$，就对应一个态射 $(\FF_q((T^{1/p^\infty})),\FF_q[\![T^{1/p^\infty}]\!])\to(R,R^+)$，由 $T\mapsto\varpi_R$ 给出。

于是投影
\[
\Bun_G\times\Spa(\FF_q((T^{1/p^\infty})))\longrightarrow\Bun_G
\]
是满态射；我们干脆宣布它是光滑的。这样问题就化归为找 $\Bun_G\times\Spa(F)$ 的一个“光滑”表示，其中 $F=\FF_q((T^{1/p^\infty}))$。下面我们从 Quot 概形的理论中汲取灵感。

我们采用如下定义。设 $S\in \mathrm{Perf}_{\FF_q}$，记 $f:(X_S)_{\text{pro-ét}}\to S_{\text{pro-ét}}$。对 $d\ge0$ 定义
\[
\mathbb{V}_{d/S}=f_*\mathcal{O}_{X_S}(d).
\]
这是 Banach–Colmez $E$-向量空间。更确切地说，若 $\mathbb{B}$ 是 $\mathrm{Perf}_{\FF_q}$ 上由
\[
\mathbb{B}(S)=\mathcal{O}(Y_S)
\]
给出的 $E$-代数层，则
\[
\mathbb{V}_d=\mathbb{B}^{\varphi=\pi^d}.
\]
当 $E=\QQ_p$ 时，这有晶体上同调的描述（更一般地，当 $E|\QQ_p$ 时用 $\pi$-除幂也可）。事实上，对完满 $\FF_q$-代数 $R$，考虑
\[
\begin{aligned}
\mathrm{B}^+_{cris}(R^\circ/\varpi_R)
&=H^0_{cris}\bigl(\Spec(R^\circ/\varpi_R)/\Spec(\ZZ_p),\mathcal{O}\bigr)\Bigl[\tfrac1p\Bigr]\\
&=W(R^\circ)\Bigl[\tfrac{\varpi_R^n}{n!}\Bigr]\Bigl[\tfrac1p\Bigr].
\end{aligned}
\]
有包含（详见 6.2 节）
\[
\mathrm{B}^+_{cris}(R^\circ/\varpi_R)\subset \mathcal{O}(Y_{R,R^+}),
\]
它诱导同构
\[
\mathrm{B}^+_{cris}(R^\circ/\varpi_R)^{\varphi=\pi^d}\xrightarrow{\ \sim\ }\mathbb{B}(R)^{\varphi=\pi^d}.
\]
这是菱形；当 $E=\FF_q((\pi))$ 时甚至是完满空间。我们有 $\mathbb{V}_0=E$。这些对象构成分次 $E$-代数
\[
\bigoplus_{d\ge0}\mathbb{V}_{d/S},
\]
所在范畴是 $S$ 上的菱形（相应地，当 $E=\FF_q((\pi))$ 时是 $S$ 上的完满空间）。

\begin{example}\label{ex:2.7}
若 $E=\FF_q((\pi))$ 且 $d>0$，考虑完满形式概形 $\mathcal{V}_d=\Spf\bigl(\FF_q[\![x_0^{1/p^\infty},\dots,x_{d-1}^{1/p^\infty}]\!]\bigr)$。它表示 1.5.2 节的形式群 $\mathcal{G}_d$ 的万有覆盖，后者是形式 $E$-向量空间。存在态射 $\mathcal{V}_{d_1}\times\mathcal{V}_{d_2}\to\mathcal{V}_{d_1+d_2}$，由展开乘积
\[
\Bigl(\sum_{i=0}^{d_1-1}\sum_{n\in\ZZ}x_i^{q^{-n}}\pi^{nd_1+i}\Bigr)\cdot\Bigl(\sum_{j=0}^{d_2-1}\sum_{m\in\ZZ}y_j^{q^{-m}}\pi^{md_2+j}\Bigr)
\]
得到。例如，$\mathcal{V}_1\times\mathcal{V}_1\to\mathcal{V}_2$ 由
\[
(x,y)\longmapsto\Bigl(\sum_{n\in\ZZ}x^{q^{-n}}y^{q^n},\ \sum_{n\in\ZZ}x^{q^{-n}}y^{q^{n-1}}\Bigr)
\]
给出。
\end{example}

\srcpage{18}

\emph{于是，取 $S$ 上的一般纤维，}
\[
\bigoplus_{d>0}\mathbb{V}_{d/S}=\bigoplus_{d>0}\mathcal{V}_d\times S,
\]
\emph{其中 $\mathcal{V}_d\times S$ 是 $d$ 维完满开球。}

菱形的任意闭子集是菱形。完满空间的任意“局部 Zariski 闭”子集是完满空间。下面的定义因此有意义。它的灵感来自：仿射代数簇由有限维向量空间中有限多个多项式的零点/非零点定义……我们于是可以对 Colmez 意义下的有限维 Banach 空间做同样的事（[4]）。

\begin{definition}\label{def:2.8}
$S$ 上的代数菱形（相应地，当 $E=\FF_q((\pi))$ 时的完满空间）是某个 $\mathbb{V}_{d/S}^n$（$d,n\ge1$）的子层，由 $E[X_1,\dots,X_n]$ 中有限多个多项式的零点与非零点定义。
\end{definition}

\begin{proposition}\label{prop:2.9}
对 $\mathcal{D}=(d_1,d_2,m_1,m_2,n)$，其中 $d_1,d_2\in\ZZ$ 满足 $d_2<d_1$，$m_1,m_2\in\mathbb{N}_{\ge1}$，且 $n\le m_1$，函子
\[
\begin{gathered}
\mathrm{Pres}_{\mathcal{D}}:\ \mathrm{Perf}_F\longrightarrow \mathrm{Sets},\\
S\longmapsto \Bigl\{u\in\Hom\bigl(\mathcal{O}_{X_S}(d_2)^{m_2},\mathcal{O}_{X_S}(d_1)^{m_1}\bigr)\ \Big|\ \coker u\ \text{局部自由、秩为 } n\Bigr\}
\end{gathered}
\]
可由 $F$ 上的代数菱形（相应地，完满空间）表示。
\end{proposition}

事实上，层
\[
S\longmapsto \Hom\bigl(\mathcal{O}_{X_S}(d_2)^{m_2},\mathcal{O}_{X_S}(d_1)^{m_1}\bigr)
\]
同构于 $\mathbb{V}^{m_1m_2}_{d_1-d_2}$，我们把它看作 $M_{m_1,m_2}(\mathbb{V}_{d_1-d_2})$，即系数在 $\mathbb{V}_{d_1-d_2}$ 中的 $m_1\times m_2$ 矩阵。现在设 $n\le m_1$。对 $u\in M_{m_1,m_2}$，$u=(X_{ij})_{1\le i\le m_1,\,1\le j\le m_2}$，考察 $E[X_{ij}]_{1\le i\le m_1,\,1\le j\le m_2}$ 中的多项式 $(P_\alpha)_\alpha$（相应地，$(Q_\beta)_\beta$），它们是 $\wedge^{n+1}u$（相应地，$\wedge^n u$）的系数。于是，所有 $(P_\alpha)_\alpha$ 的零点与 $(Q_\beta)_\beta$ 中某一个的非零点，在 $\mathbb{V}^{m_1m_2}_{d_1-d_2}$ 中定义了一个代数菱形。它表示 $\mathrm{Perf}_{\mathcal{D}}$。

有明显的态射
\[
\mathrm{pres}_{\mathcal{D}}:\ \mathrm{Pres}_{\mathcal{D}}\longrightarrow \Bun_G\times_{\Spa(\FF_q)}\Spa(F).
\]
于是可以验证下面的事实。

\begin{proposition}\label{prop:2.10}
我们可以把 $\Bun_G$ 写成拟紧开子集的可数并
\[
\Bun_G=\bigcup_{n\ge1}U_n,
\]
使得对每个 $n$，存在如命题 2.9 中的数据 $\mathcal{D}_n$，满足：
\begin{enumerate}[label=(\arabic*),leftmargin=*]
\item $U_n\times\Spa(F)\subset \mathrm{Im}(\mathrm{pres}_{\mathcal{D}_n})$；
\item 限制在 $\mathrm{pres}_{\mathcal{D}_n}^{-1}(U_n\times\Spa(F))$ 上的 $\mathrm{pres}_{\mathcal{D}_n}$ 的几何纤维满足：
\begin{itemize}
\item 若 $E=\FF_q((\pi))$，它们同构于维数相同的完满开球（与 $U_n$ 的几何点无关）；
\item 若 $E|\QQ_p$，存在 $d\ge1$ 与 $h\in\mathbb{N}$，使得对每个几何纤维 $Z$，存在 pro-étale Galois 覆盖 $\widetilde{Z}\to Z$，其 Galois 群为 $E^h$，且 $\widetilde{Z}$ 同构于 $d$ 维完满开球。
\end{itemize}
\end{enumerate}
\end{proposition}

\yizhuw{原文此处写“Galois 群 $E^h$”。$E$ 是 $p$-进域，而 $E^h$ 不是常见的 Galois 群记号，疑为 $(\ZZ/p^h)$ 或 $\FF_p^h$ 一类；译文照原文保留，未改。}

无论怎样理解这句话，态射
\[
\bigl(\mathrm{pres}_{\mathcal{D}_n}\bigr)\big|_{\mathrm{pres}^{-1}_{\mathcal{D}_n}(U_n)}
\]
都是光滑且满的。这给了我们 $\Bun_G$ 的一个“表示”。\srcpage{19}

\subsection{叠的点}% 原文 2.2. Points of the stack

固定 $\FF_q$ 的一个代数闭包 $\ol{\FF}_q$，并记 $L=\widehat{E^{un}}$ 为 $E$ 的相应极大非分歧扩域的完备化。设 $\sigma$ 是它的 Frobenius。回忆（Kottwitz [18]）：
\[
B(G):=G(L)/\sigma\text{-共轭}.
\]
我们记
\[
\varphi\text{-}\mathrm{Mod}_L=\{(D,\varphi)\}
\]
为相应的等晶体范畴，其分类由 Dieudonné 与 Manin 给出；这里 $D$ 是有限维 $L$-向量空间，$\varphi$ 是 $D$ 的 $\sigma$-线性自同构。每个 $b\in G(L)$ 对应一个带 $G$-结构的等晶体
\[
\begin{array}{rcl}
\mathcal{F}_b:\ \Rep(G) & \longrightarrow & \varphi\text{-}\mathrm{Mod}_L\\[2pt]
(V,\rho) & \longmapsto & (V\otimes_E L,\rho(b)\sigma).
\end{array}
\]
由此得到一个双射
\[
\begin{array}{rcl}
\text{$G$-等晶体}/\sim & \xrightarrow{\ \sim\ } & B(G)\\[2pt]
[b] & \longmapsto & [\mathcal{F}_b].
\end{array}
\]

\yizhu{$B(G)$ 即 Kottwitz 集：$G(L)$ 模 $\sigma$-共轭（$b\mapsto gb\sigma(g)^{-1}$）所得的商集。}{A：原文 §2.2 p.19；Kottwitz [18]}

\yizhuw{原文的对应“$[b]\mapsto[\mathcal{F}_b]$”与上一行箭头方向（$G$-等晶体 $\to B(G)$）不相合，似应读作 $[\mathcal{F}_b]\mapsto[b]$；译文照原文排，未改。}

\subsubsection*{2.2.1. 向量丛的分类}% 原文 2.2.1. Classification of vector bundles

现在固定 $\ol{\FF}_q$ 上的完满域 $F$，考虑相应的曲线 $X^{ad}$ 或 $X^{sch}$。有函子
\[
\begin{array}{rcl}
\varphi\text{-}\mathrm{Mod}_L & \longrightarrow & \Bun_{X^{ad}}\\[2pt]
(D,\varphi) & \longmapsto & \mathcal{E}(D,\varphi)^{ad},
\end{array}
\]
其定义如下。由于 $F|\ol{\FF}_q$，空间 $Y$ 落在 $\Spa(L)$ 之上。$\mathcal{E}(D,\varphi)^{ad}$ 的几何实现由
\[
Y\times_\varphi D\longrightarrow Y/\varphi^{\ZZ}=X^{ad}
\]
给出。也就是说，向量丛 $\mathcal{E}(D,\varphi)^{ad}$ 沿覆盖 $Y\to X^{ad}$ 拉回到 $Y$ 后，变成以 $D$ 为纤维的平凡丛；相应的自守因子由 $\varphi$ 在 $D$ 上的作用给出。由 GAGA
\[
\Bun_{X^{sch}}\xrightarrow{\ \sim\ }\Bun_{X^{ad}},
\]
向量丛 $\mathcal{E}(D,\varphi)^{ad}$ 是代数向量丛 $\mathcal{E}(D,\varphi)$ 的解析化，后者对应分次代数
\[
\bigoplus_{d\ge0}\bigl(D\otimes_L\mathcal{O}(Y)\bigr)^{\varphi\otimes\varphi=\pi^d},
\]
它定义在 $P=\bigoplus_{d\ge0}\mathcal{O}(Y)^{\varphi=\pi^d}$ 之上，且 $X^{sch}=\Proj(P)$。于是有如下定理。

\begin{theorem}\label{thm:2.11}
（[10]）若 $F$ 代数闭，则函子 $\mathcal{E}(-):\varphi\text{-}\mathrm{Mod}_L\to \Bun_X$ 本质满。
\end{theorem}

借助 Dieudonné–Manin 分类，这个定理可以有更具体的表述。事实上，对每个 $d\in\ZZ$，我们已经定义了 $X$ 上的线丛 $\mathcal{O}(d)$。有 $\mathcal{O}(d)=\mathcal{E}(L,\pi^{-d}\sigma)$。对整数 $h\ge1$，记 $E_h$ 为 $E$ 的 $h$ 次非分歧扩域。有
\[
X_E\otimes_E E_h=X_{E_h},
\]
这个公式反映了：把 $E$ 换成 $E_h$ 不改变 Fontaine 环 $\mathbf{A}$，但把 $\varphi$ 换成 $\varphi^h$。循环 Galois 覆盖
\[
X^{ad}_{E_h}\longrightarrow X^{ad}_E
\]
于是等同于展开覆盖
\[
Y/\varphi^{h\ZZ}\longrightarrow Y/\varphi^{\ZZ}.
\]
我们记 $X_h:=X_{E_h}$，$X_1=X$，并有 $\pi_h:X_h\to X$。于是有典范的 $\widehat{\ZZ}$-Galois 覆盖
\[
(X_h)_{h\ge1}\longrightarrow X.
\]

\yizhu{“展开覆盖”指 $X^{ad}_{E_h}\to X^{ad}_E$ 等同于 $Y/\varphi^{h\ZZ}\to Y/\varphi^{\ZZ}$：曲线 $X$ 是 $Y$ 模 $\varphi$-作用的商，换成 $h$ 次非分歧扩域就是把作用换成 $\varphi^h$。}{A：原文 §2.2.1 p.19}

\srcpage{20}

对 $\lambda\in\QQ$，写 $\lambda=d/h$ 且 $(d,h)=1$，定义
\[
\mathcal{O}(\lambda):=\pi_{h*}\mathcal{O}_{X_h}(d).
\]
有 $\mathcal{O}(\lambda)\simeq\mathcal{E}(D,\varphi)$，其中 $(D,\varphi)$ 是单的，其 Dieudonné–Manin 斜率为 $-\lambda$。于是前述定理可以改述为：\emph{当 $F$ 代数闭时，对任意 $\mathcal{E}\in\Bun_X$，存在有限族斜率 $(\lambda_i)_{i\in I}$，使得}
\[
\mathcal{E}\simeq\bigoplus_{i\in I}\mathcal{O}(\lambda_i).
\]

还可以再往前走一步地改述这个定理。事实上，曲线 $X^{sch}$ 是完备的，意即对任意 $f\in E(X^{sch})^\times$，有 $\deg(\mathrm{div}\,f)=0$。因此 $X^{sch}$ 上的向量丛有良好的次数函数，向量丛上也有 Harder–Narasimhan 滤过。对每个 $\lambda\in\QQ$，$\mathcal{O}(\lambda)$ 是斜率为 $\lambda$ 的稳定丛。于是前述定理可以改述如下。

\begin{theorem}\label{thm:2.12}
若 $F$ 代数闭：
\begin{enumerate}[label=(\arabic*),leftmargin=*]
\item 对任意 $\lambda\in\QQ$，任意斜率为 $\lambda$ 的半稳定向量丛都同构于 $\mathcal{O}(\lambda)$ 的有限直和。
\item $X^{sch}$ 上向量丛的 Harder–Narasimhan 滤过是分裂的。
\end{enumerate}
\end{theorem}

我们指出，第二点 (2) 对任意 $F$ 仍成立，不必代数闭。

于是有双射
\[
\begin{array}{rcl}
\{\lambda_1\ge\cdots\ge\lambda_n\ |\ n\in\mathbb{N},\ \lambda_i\in\QQ\} & \xrightarrow{\ \sim\ } & \Bun_X/\sim\\[2pt]
(\lambda_1,\dots,\lambda_n) & \longmapsto & \Bigl[\bigoplus_{i=1}^{n}\mathcal{O}(\lambda_i)\Bigr].
\end{array}
\]

\subsubsection*{2.2.2. $G$-丛的分类}% 原文 2.2.2. Classification of G-bundles

设 $b\in G(L)$。复合函子
\[
\Rep(G)\xrightarrow{\ \mathcal{F}_b\ }\varphi\text{-}\mathrm{Mod}_L\xrightarrow{\ \mathcal{E}(-)\ }\Bun_X
\]
定义了 $X$ 上一个带 $G$-结构的丛，也就是说概形 $X^{sch}$ 上一个对 étale 拓扑局部平凡的 $G$-torsor
\[
\mathcal{E}_b.
\]
于是得到向量丛分类定理的如下推广。

\begin{theorem}\label{thm:2.13}
若 $E|\QQ_p$ 且 $F$ 代数闭，则存在带点集的双射
\[
\begin{array}{rcl}
B(G) & \xrightarrow{\ \sim\ } & H^1_{\text{ét}}(X^{sch},G)\\[2pt]
[b] & \longmapsto & [\mathcal{E}_b].
\end{array}
\]
\end{theorem}

当然，当 $E=\FF_q((\pi))$ 时这也必须成立。

\begin{conjecture}\label{conj:2.14}
当 $E=\FF_q((\pi))$ 时，定理 2.13 仍然成立。
\end{conjecture}

我们将承认这个猜想，并在本文余下部分把它当作真的来用。

当 $G$ 拟分裂时，这个双射和前面一样带有 Harder–Narasimhan 的特征。为此固定三元组 $A\subset T\subset B$，其中 $A$ 是极大分裂环面，$T$ 是极大环面，$B$ 是 Borel 子群。Kottwitz 定义了斜率映射
\[
\begin{array}{rcl}
B(G) & \longrightarrow & \bigl[\Hom(\mathbb{D},G_{E^{un}})/\mathrm{conj.}\bigr]^{\sigma=Id}=X_*(A)^+_{\QQ}\\[2pt]
[b] & \longmapsto & [\nu_b].
\end{array}
\]
曲线 $X^{sch}$ 上的 $G$-torsor 有半稳定性概念，它与“经典曲线”上通常的半稳定性类似，用约化到 $G$ 的标准抛物子群来表述。按定义，若 $\nu_b$ 是中心的，则 $b$ 是基本的。于是有如下等价：
\[
b\ \text{是基本的}\iff \mathcal{E}_b\ \text{是半稳定的}.
\]

\srcpage{21}

此外，若 $w$ 是 $T$ 的 Weyl 群中的最长元，使得 $B^w$ 与 $B$ 相反，则
\[
[w.(-\nu_b)]\in X_*(A)^+_{\QQ}
\]
是 $\mathcal{E}_b$ 的广义 Harder–Narasimhan 多边形。这里的负号来自：在对应 $(D,\varphi)\mapsto\mathcal{E}(D,\varphi)$ 之下，Dieudonné–Manin 斜率与 Harder–Narasimhan 斜率相反。

\subsubsection*{2.2.3. 叠的点}% 原文 2.2.3. Points of the stack

任意 $b\in G(L)$ 定义态射
\[
x_b:\ \Spa(\ol{\FF}_q)\longrightarrow \Bun_G.
\]
我们置
\[
\Bun_{G,\ol{\FF}_q}:=\Bun_G\otimes_{\FF_q}\ol{\FF}_q.
\]
现在定义
\[
\bigl|\Bun_{G,\ol{\FF}_q}\bigr|=\Bigl(\coprod_{F}\Bun_G(F)\Bigr)\Big/\sim,
\]
其中 $F$ 遍历 $\ol{\FF}_q$ 的一切完满扩域。按定义，对 $x_1\in\Bun_G(F_1)$、$x_2\in\Bun_G(F_2)$，若存在 $F_3$ 以及嵌入 $\alpha:F_1\hookrightarrow F_3$、$\beta:F_2\hookrightarrow F_3$，使得 $\alpha_*x\simeq\beta_*y$，则 $x_1\sim x_2$。

由前述分类定理可得，存在典范双射
\[
B(G)\xrightarrow{\ \sim\ }\bigl|\Bun_{G,\ol{\FF}_q}\bigr|.
\]
现在给 $\bigl|\Bun_{G,\ol{\FF}_q}\bigr|$、从而给 $B(G)$ 装上这样的拓扑：其开子集是 $|\mathcal{U}|$，其中 $\mathcal{U}\subset\Bun_{G,\ol{\FF}_q}$ 是开子叠。要注意，这与公式 $B(G)=G(L)/\sigma\text{-conj.}$ 给出的商拓扑不同（后者是离散拓扑）。

\yizhuw{原文在“$x_1\sim x_2$”的定义里用了 $\alpha_*x\simeq\beta_*y$，而前文的对象记作 $x_1,x_2$；$x,y$ 当是笔误，译文照原文保留。}

\subsection{连通分支}% 原文 2.3. Connected components

Kottwitz 定义了映射
\[
\kappa:\ B(G)\longrightarrow \pi_1(G)_\Gamma,
\]
其中 $\Gamma=\Gal(\ol{E}|E)$，$\pi_1(G)$ 是 Borovoi 的基本群。当 $G=\mathrm{GL}_n$ 时，这给出 Newton 多边形的端点，也就是相应向量丛次数的相反数。一般地，这可以解释为一个 $G$-torsor 的 $G$-等变第一 Chern 类（[6] 第 8 节）。下述定理在 $G_{der}$ 单连通时是容易的（此时经由到 $G/G_{der}$ 的投影化归到环面情形），但一般情形微妙得多。

\begin{theorem}\label{thm:2.15}
Kottwitz 映射 $\kappa$ 在 $\Bun_{G,\ol{\FF}_q}$ 上局部常值。
\end{theorem}

于是有开/闭子叠的分解
\[
\Bun_{G,\ol{\FF}_q}=\coprod_{\alpha\in\pi_1(G)_\Gamma}\Bun^{\alpha}_{G,\ol{\FF}_q}.
\]
下述猜想与“经典情形”类比，显得自然。陈述我们的猜想时并不需要它。

\begin{conjecture}\label{conj:2.16}
对每个 $\alpha\in\pi_1(G)_\Gamma$，$\Bun^{\alpha}_{G,\ol{\FF}_q}$ 连通。
\end{conjecture}

\subsection{\texorpdfstring{$\Bun_G$}{Bun\_G} 的 Harder–Narasimhan 分层}% 原文 2.4. HN stratification of Bun_G

\emph{从现在起我们假设 $G$ 拟分裂。}

利用 Kedlaya 与 Liu 的一个结果（他们处理了 $\mathrm{GL}_n$ 情形，见 [16]），我们可以证明下面的事实。我们给 $X_*(A)^+_{\QQ}$ 装上这样的序：若 $\nu_2-\nu_1$ 是余根的正 $\QQ$-线性组合，则 $\nu_1\le\nu_2$。

\begin{theorem}\label{thm:2.17}
Harder–Narasimhan 映射 $\HN:B(G)\to X_*(A)^+_{\QQ}$，$[b]\mapsto[w.(-\nu_b)]$，是半连续的，意即对一切 $\nu\in X_*(A)^+_{\QQ}$，使 $\HN([b])\le\nu$ 的 $[b]\in B(G)$ 构成闭集。
\end{theorem}

\srcpage{22}

陈述我们的猜想时，只需用到下面这个推论。

\begin{corollary}\label{cor:2.18}
半稳定轨迹 $\Bun_G^{ss}$ 是开子叠。
\end{corollary}

\subsection{\texorpdfstring{$\mathcal{E}_b$}{E\_b} 的自同构群}% 原文 2.5. The automorphism group of E_b

固定 $b\in G(L)$，考虑 $b$ 的 $\sigma$-中心化子 $J_b$。它是 $G$ 的某个 Levi 子群的内形式，即斜率态射 $\nu_b$ 的中心化子。特别地，若 $b$ 是基本的，它就是 $G$ 的内形式。

\begin{definition}\label{def:2.19}
我们记 $\wt{J}_b$ 为 $\mathrm{Perf}_{\ol{\FF}_q}$ 上由
\[
\wt{J}_b(S)=\Aut(\mathcal{E}_b|_{X_S})
\]
定义的 pro-étale 层。
\end{definition}

由 2.6，这是一个菱形群（更精确地说，限制到任意完满域 $F$ 上的 $\mathrm{Perf}_F$ 时如此）。有
\[
J_b(E)=\Aut(\mathcal{F}_b),
\]
从而有包含
\[
\ul{J_b(E)}\subset \wt{J}_b.
\]
于是有如下命题。

\begin{proposition}\label{prop:2.20}
若 $b$ 是基本的，则 $\ul{J_b(E)}=\wt{J}_b$。
\end{proposition}

\begin{example}\label{ex:2.21}
平凡 $G$-丛 $\mathcal{E}_1$ 的自同构由 $G(E)$ 给出。更确切地说，对任意 $S\in \mathrm{Perf}_{\ol{\FF}_q}$，$X_S$ 上平凡 $G$-丛的自同构群是连续映射构成的群 $\mathcal{C}(|S|,G(E))$。
\end{example}

\begin{remark}\label{rem:2.22}
与“经典情形”（代数曲线上的 $G$-丛叠，$G$ 定义在域 $k$ 上）相反，我们看到平凡 $G$-丛的自同构是 $G(E)$，而不是代数群 $G$！
\end{remark}

但 $\mathcal{E}_b$ 的任意自同构诱导其典范半稳定约化的自同构，因此总有态射
\[
\wt{J}_b\longrightarrow J_b(E).
\]
最后，关于 $\wt{J}_b$ 有如下结构结果。

\begin{proposition}\label{prop:2.23}
有 $\pi_0(\wt{J}_b)=J_b(E)$，且 $\wt{J}_b=\wt{J}_b^0\rtimes\ul{J_b(E)}$。此外，单位连通分支 $\wt{J}_b^0$ 是幂单群菱形；当 $E=\FF_q((\pi))$ 时，是完满空间。
\end{proposition}

\begin{example}\label{ex:2.24}
看一个简单的例子。取 $G=\mathrm{GL}_2$，$b=\mathrm{diag}(1,\pi^{-1})$。此时 $\mathcal{E}_b=\mathcal{O}\oplus\mathcal{O}(1)$。于是
\[
\wt{J}_b=\begin{pmatrix}\ul{E^\times} & \mathbb{B}^{\varphi=\pi}\\ 0 & \ul{E^\times}\end{pmatrix}.
\]
由此可见，$\wt{J}_b^0=\mathbb{B}^{\varphi=\pi}$，是 Banach–Colmez 空间。
\end{example}

一般地，$\wt{J}_b^0$ 有滤过，其分次片是 Banach–Colmez 空间。

\subsection{HN 分层的单值化}% 原文 2.6. Uniformization of HN strata

\subsubsection*{2.6.1. 基本情形}% 原文 2.6.1. The basic case

我们先处理基本情形，因为它更简单，而且陈述我们的猜想只需要这一情形。回忆（Kottwitz）：$\kappa$ 的限制诱导双射
\begin{equation}\tag{7}\label{eq:2.7}
\kappa:\ B(G)_{basic}\xrightarrow{\ \sim\ }\pi_1(G)_\Gamma.
\end{equation}
几何上这翻译为如下说法。

\begin{proposition}\label{prop:2.25}
对每个 $\alpha\in\pi_1(G)_\Gamma$，由 $\alpha$ 切出的分支中的基本轨迹，即开子集
\[
\bigl|\Bun^{\alpha,ss}_{G,\ol{\FF}_q}\bigr|,
\]
只是由 (7) 决定的一个点。
\end{proposition}

\srcpage{23}

固定 $b\in G(L)$，使 $\kappa([b])=\alpha$。于是有态射
\begin{equation}\tag{8}\label{eq:2.8}
x_b:\ \bigl[\Spa(\ol{\FF}_q)/\ul{J_b(E)}\bigr]\longrightarrow \Bun^{\alpha,ss}_{G,\ol{\FF}_q},
\end{equation}
由 $\mathcal{E}_b$ 定义。这里左端是 $J_b(E)$-pro-étale torsor 的分类叠。由 Kedlaya 与 Liu 的结果（$b=1$ 的 $\mathrm{GL}_n$ 情形，见 [16]），可以证明下面的事实。

\yizhu{torsor 即主齐性空间（带群自由传递作用的对象），此处保留英文。}{B：Stacks Project，“Torsors”}

\begin{theorem}\label{thm:2.26}
态射 (8) 是同构
\[
\bigl[\Spa(\ol{\FF}_q)/\ul{J_b(E)}\bigr]\xrightarrow{\ \sim\ }\Bun^{\alpha,ss}_{G,\ol{\FF}_q}.
\]
\end{theorem}

这个定理的意思只是：若 $S\in \mathrm{Perf}_{\ol{\FF}_q}$，$\mathcal{E}$ 是 $X_S$ 上的 $G$-丛，且对每个几何点 $\ol{s}\to S$ 有 $\mathcal{E}|_{X_{k(\ol{s})}}\simeq\mathcal{E}_b$，则存在 pro-étale 覆盖 $T\to S$，使得 $\mathcal{E}|_{X_T}\simeq\mathcal{E}_b$。相应的 $S$ 上的 $\ul{J_b(E)}$-torsor，就是 $\mathcal{E}_b$ 与 $\mathcal{E}$ 之间同构构成的 torsor。

\subsubsection*{2.6.2. 一般情形}% 原文 2.6.2. The general case

陈述猜想时用不到这个一般情形，读者可以跳过。

据 Kottwitz，映射
\[
B(G)\xrightarrow{\ (\kappa,\nu)\ }\pi_1(G)_\Gamma\times X_*(A)^+_{\QQ}
\]
是单射。几何上这翻译为如下说法。

\begin{proposition}\label{prop:2.27}
固定 $\alpha\in\pi_1(G)_\Gamma$。则对每个 $\nu\in X_*(A)^+_{\QQ}$，分层
\[
\bigl|\Bun^{\alpha,HN=\nu}_{G,\ol{\FF}_q}\bigr|
\]
或者为空，或者是一个点。
\end{proposition}

设这样的分层非空，并固定 $b$，使得 $\kappa(b)=\alpha$、$w.(-\nu_b)=\nu$。则定理 2.26 在此情形下推广为
\[
\bigl[\Spa(\ol{\FF}_q)/\wt{J}_b\bigr]\xrightarrow{\ \sim\ }\Bun^{\alpha,HN=\nu}_{G,\ol{\FF}_q},
\]
其中右端是分类这样的 $G$-丛的叠：它们在几何纤维上同构于 $\mathcal{E}_b$。

无论按哪种能想象的意义，都有
\[
\dim\wt{J}_b=\langle 2\rho,\mu\rangle,
\]
其中 $\rho$ 是 $T$ 在 $\Lie B$ 中正权之和的一半。于是
\[
\dim\Bun^{\alpha,HN=\nu}_{G,\ol{\FF}_q}=-\langle 2\rho,\mu\rangle.
\]
由此得到：
\begin{itemize}
\item \emph{基本轨迹是零维的}；
\item 当我们深入 Weyl 房时，HN 分层的维数趋于 $-\infty$。
\end{itemize}
